\section{Discusión}
~\label{sec:discusion}

Los resultados presentados en la Sección~\ref{sec:resultados} 
confirman que el sistema opera dentro de los márgenes previstos 
por el diseño y satisface los requerimientos funcionales y no 
funcionales para los que fue construido. Sin embargo, una lectura 
útil del trabajo exige ir más allá de la verificación de 
cumplimiento: obliga a examinar las tensiones que atravesaron las 
decisiones arquitectónicas, los márgenes entre lo que el sistema 
logra y lo que queda fuera de su alcance, y la posición del 
proyecto dentro del espacio más amplio del monitoreo de red en 
redes \@. Esta sección organiza esa reflexión en cuatro ejes.

\subsection{Tensiones arquitectónicas y costos asumidos} 

La captura distribuida a nivel de endpoint, que resuelve la 
atribución de comportamiento por dispositivo identificada como 
causa raíz C2, asume un costo operativo en donde cada endpoint debe alojar 
un proceso adicional, el administrador pierde parte del control 
centralizado característico de los modelos basados en port mirror, 
y la fase de despliegue se multiplica por el número de 
dispositivos. El modelo centralizado clásico evita ese costo a 
cambio de perder precisamente la granularidad que justifica este 
proyecto. La elección del modelo distribuido es coherente con el 
objetivo general, pero el costo de despliegue no es trivial y se 
vuelve más significativo conforme crece el tamaño de la red.

La sanitización dinámica reduce el volumen transmitido y protege 
la confidencialidad del payload, pero descarta información que 
podría ser diagnóstica en escenarios específicos. El recorte a 
96 o 300 bytes preserva las cabeceras y los campos diagnósticos 
del nivel de aplicación para DNS y TLS, pero inhabilita cualquier 
análisis que requiera inspección del payload completo, como la 
detección de patrones de exfiltración codificados en HTTP o la 
identificación de malware mediante firmas en el contenido. El 
diseño asume este costo de forma consciente porque el proyecto se 
orienta al monitoreo de comportamiento agregado, no a la 
detección basada en firmas; sin embargo, al pasar del laboratorio 
a un despliegue de producción, este límite debe evaluarse frente 
a los escenarios específicos que el operador necesite cubrir.

La separación del pipeline y la API como procesos independientes 
comunicados por Shared Database aísla fallos y permite 
escalabilidad independiente, pero introduce latencia adicional 
entre la generación de una métrica y su disponibilidad para la 
consulta móvil: el dato atraviesa la escritura a PostgreSQL, el 
mecanismo \texttt{LISTEN/NOTIFY} y el handler WebSocket antes de 
llegar al cliente. La alternativa de acoplamiento directo por 
memoria compartida sería más rápida, pero propagaría fallos del 
pipeline a la capa de exposición, con consecuencias más severas 
para la disponibilidad del sistema desde la perspectiva del 
administrador. Para un sistema de monitoreo donde un retardo de 
segundos es aceptable pero una caída no lo es, el intercambio 
parece correcto; para un sistema de respuesta en tiempo real, 
probablemente no lo sería.

\subsection{Detección de anomalías: poder y límites}

El motor basado en EWMA + Z-Score adaptativo con filtro de 
persistencia resuelve limitaciones concretas de los enfoques más 
simples. Los umbrales fijos no distinguen el perfil individual de 
cada dispositivo y fracasan simultáneamente en dos direcciones 
opuestas: generan falsas alarmas en dispositivos que operan 
legítimamente cerca del umbral, y pierden anomalías en 
dispositivos cuyo comportamiento normal está lejos del umbral. El 
Z-Score con media acumulada incorpora el perfil individual pero 
tarda horas en adaptarse a cambios legítimos de régimen, lo que 
produce avalanchas de falsas alarmas durante los períodos de 
convergencia. El método adoptado resuelve ambas limitaciones a 
cambio de una complejidad modesta: dos variables de estado por 
par dispositivo-métrica, complejidad temporal constante por 
ciclo, y un ajuste único del factor de suavizado.

La discusión anterior, sin embargo, se queda corta si no aborda lo 
que el motor no puede detectar. Una anomalía que respete los 
rangos individuales de cada métrica pero cuya combinación sea 
inusual permanecerá invisible para un detector que evalúa cada 
métrica de forma independiente. Por ejemplo, un incremento 
moderado de \texttt{bandwidth\_out} junto con un incremento 
moderado de \texttt{unique\_destinations} y una caída simultánea 
de \texttt{dns\_response\_time} es un patrón que, tomadas las 
métricas por separado, no superaría el umbral de severidad, pero 
colectivamente podría indicar actividad de escaneo previo a un 
ataque. Detectar este tipo de correlación requiere un enfoque 
multivariado que excede el alcance del motor implementado.

Análogamente, las anomalías de secuencia, donde ningún valor 
individual es anómalo pero sí lo es el patrón temporal, tampoco 
son capturadas por el detector actual. Un tráfico sostenido durante 
seis horas fuera de horario laboral puede estar compuesto por 
valores individualmente normales que el baseline EWMA absorbe como 
la nueva línea base; la anomalía reside en el contexto temporal 
que el sistema actual no modela. Chandola et al.~\cite{chandola2009anomaly} 
denominan a estas tres categorías anomalías puntuales, 
contextuales y colectivas; el motor cubre las dos primeras y 
declara abierta la tercera como trabajo futuro.

\subsection{Validación frente a objetivos y contribución 
disciplinar}

El proyecto satisface sus cuatro objetivos específicos en un 
grado variable que conviene explicitar. El diseño arquitectónico 
(OE1) produjo una solución coherente con el estado del arte del 
monitoreo distribuido: la correspondencia entre los sensores y los 
\textit{observation points} de la RFC 5470, entre el backend y el 
\textit{collector} de la misma RFC\@, y entre el flujo interno y el 
patrón \textit{Pipes and Filters}~\cite{buschmann1996posa} sitúa 
el trabajo sobre referentes establecidos en lugar de inventar una 
organización ad hoc. La implementación (OE2) produjo artefactos 
ejecutables para los tres componentes principales, respaldados por 
una batería de pruebas automatizadas que valida las rutas críticas 
del sistema. La validación funcional (OE3) se cubrió mediante la 
ejecución de escenarios controlados con detección efectiva de 
anomalías inducidas. El despliegue mediante pipelines de 
integración continua (OE4) quedó operativo a nivel de backend y 
de la imagen móvil, si bien la publicación del APK firmado en 
tiendas formales no se cubrió por estar fuera del alcance 
académico del proyecto.

La contribución del trabajo no reside en la originalidad de 
ninguna técnica aislada. EWMA data de 1959~\cite{roberts1959ewma}, 
el patrón Thread-per-Client está documentado desde hace 
décadas~\cite{schmidt2000posa2}, y la arquitectura 
sensor-colector se formalizó en la RFC 5470 hace quince años. La 
contribución reside, en cambio, en la integración coherente de 
estos elementos en un sistema funcional acoplado a una aplicación 
móvil, bajo las restricciones de un proyecto integrador académico 
y con una documentación rastreable de cada decisión. El valor del 
trabajo es la demostración de que los componentes disponibles del 
estado del arte, combinados con criterio de ingeniería, producen 
una solución operativa para un problema que en la práctica 
docente suele abordarse con herramientas comerciales de costo 
considerable.

\subsection{Gap entre laboratorio y producción}

Es útil cerrar la discusión reconociendo la distancia entre un 
sistema que funciona en el entorno de validación descrito en la 
Sección~\ref{sec:resultados} y un sistema listo para un despliegue 
operativo real en una organización. Tres aspectos marcan esa 
distancia.

El primero es la sensibilidad al ajuste de parámetros. El factor 
de suavizado del EWMA, los umbrales de severidad, los parámetros 
del filtro de persistencia y los valores mínimos de varianza son 
hiperparámetros que en el presente trabajo se fijaron mediante 
justificación teórica y validación empírica sobre escenarios 
acotados. Un despliegue en producción exigiría una fase de ajuste 
específico al perfil de tráfico de la red a monitorear, 
posiblemente automatizado mediante optimización sobre datos 
históricos, que el proyecto no abordó.

El segundo es la superficie de ataque del propio sistema. El 
sensor de captura ejecuta con privilegios elevados para acceder a 
la interfaz de red en modo promiscuo, el backend ingiere datos 
binarios arbitrarios de la red, y la API expone endpoints 
autenticados por JWT\@. Cada una de estas puertas es un vector 
potencial de ataque que requeriría, para un despliegue real, 
análisis de seguridad especializado, endurecimiento de las 
configuraciones por defecto, auditoría del código para 
vulnerabilidades conocidas y quizá un programa de revelación 
coordinada de vulnerabilidades. El presente trabajo incorporó 
medidas defensivas razonables (aislación de Zeek en contenedor, 
firma criptográfica de tokens, cifrado del canal), pero no se 
sometió a una evaluación de seguridad formal.

El tercero es la operabilidad a largo plazo. Un sistema en 
producción requiere monitoreo de sí mismo (observabilidad), 
procedimientos documentados de recuperación ante fallos, 
políticas de retención y respaldo de datos, y capacidad de 
actualización sin interrupciones. El proyecto produjo un sistema 
funcional pero no un producto operable en ese sentido más amplio. 
La distancia entre un prototipo académico validado y un producto 
desplegable no es una objeción al trabajo, sino una descripción 
realista de su alcance.